\section{Podsumowanie}
\label{sec:podsumowanie}
W ramach pracy stworzono stanowisko testowe do przeprowadzania badań HIL układu wykonawczego systemu sterowania rakiety sterowanej FOK 2.
Przeprowadzono wstępne badania używając stanowiska testowego, tym samym potwierdzając jego poprawne działanie
i użyteczność w badaniach HIL. Przeprowadzone badania uwidoczniły też kilka aspektów, które wymagają poprawy w działaniu stanowiska.
W szczególności zidentyfikowano potrzebę wprowadzenia korekty pomiaru wychyleń sterów w oprogramowaniu stanowiska 
uwidoczniono wrażliwość stanowiska na względne ustawienie czujnika i magnesu. 
Zidentyfikowano obszary, w których należy w najbliższym czasie ulepszyć stanowisko testowe.

Ponadto opisano proces przeprowadzania symulacji HIL na komputerze osobistym z użyciem \textit{Simulink}a
oraz na komputerze czasu rzeczywistego \textit{Speedgoat}. Pozwoli to członkom
Zespołu na samodzielne przeprowadzanie symulacji HIL w przyszłości i ułatwi proces przekazywania wiedzy o działaniu 
symulacji HIL.

Stworzone stanowisko testowe jest integralną częścią stanowiska 
testowego całej rakiety FOK 2, które umożliwia przeprowadzanie badań HIL pętli sterowania zawierającej
zarówno głowicę obrazującą, czujnik INS/GNSS, komputer pokładowy i układ wykonawczy systemu sterowania rakiety.
Nabyte możliwości w zakresie przeprowadzania symulacji HIL pozwolą na lepszą weryfikację
działania rakiety FOK 2 przed planowanymi testami lotnymi. W szczególności 
na sprawdzenie poprawności działania pętli sterowania rakiety, jej oprogramowania pokładowego oraz rozważenie scenariuszy awaryjnych, które mogą wystąpić podczas testów lotnych. 